\chapter{시간: 언제 무엇을 하는가}\label{ch:relay-timing}

실물 기계에서 이어진 선을 찾는 일은 전류가 공짜로 하고, 남는 일은 시간을 맞추는 것이다(\ref{sec:f-current}절).
연결이 자리 잡은 뒤에 전압을 걸고, 감지 릴레이가 붙을 만큼 기다려 판정하고, 연결이 바뀌기 전에 지워야 한다.
어려운 점은 부품이다. 릴레이의 지연은 나노초가 아니라 밀리초이고, 같은 품번이라도 개체마다 두 배씩 다르다.
이 장은 그런 부품으로 틀림없이 동작하는 순서를 짜는 방법을 다룬다. 테스트 레지스터의 판정 순서, 그 순서를
지키게 하는 시간 규칙, 그리고 두 기계가 함께 쓰는 커넥터의 약속이다. 규칙을 실제로 지키는 장치는 기계마다
다르다. 기어식은 슬롯 원판(\ref{sec:t-gear}절), 릴레이식은 순서기(\ref{ch:relay-sequencer}장)이고, 각각 D부와
E부에서 본다. FPGA에서는 설계
도구가 모든 경로의 지연을 계산해 같은 종류의 규칙을 확인해 준다(\ref{sec:timing}절).

\section{들쭉날쭉한 부품으로 시간을 정하기}\label{sec:t-bounds}

\Why 이 프로젝트가 쓰는 릴레이 시간은 표~\ref{tab:t-relay}와 같다. 데이터시트에 적힌 것은 대부분 ``최대''이고,
최소는 알 수 없는 경우가 많다. 그러면 시간 규칙은 모두 \emph{하한}으로 적어야 한다. ``\code{STROBE}는 7.5~ms
\emph{이상}'', ``\code{JUDGE}는 \code{STROBE}가 오른 뒤 11.5~ms \emph{이후}''처럼. 하한으로 적은 규칙에는 좋은
성질이 있다. 기계가 \emph{느려지면} 언제나 더 잘 지켜진다. 손으로 천천히 돌리든, 관람객이 한 걸음씩
누르든, 모터가 가속하는 중이든 상관없다. 위험한 것은 기계의 박자가 빨라지는 쪽 하나뿐이다. 이 장의 설계 결정은 대부분 이
성질에서 나온다.

\begin{table}[ht]
\centering
\caption{이 프로젝트가 쓰는 릴레이 시간(결정 0001). 일부는 데이터시트의 대표·최대값이고 일부는 추정이다.
시제품에서 실측해 확정한다.}
\label{tab:t-relay}
\small
\begin{tabular}{lrr}
\toprule
& 대표 & 범위 \\
\midrule
신호 릴레이가 붙는 시간(접점 안정 포함) & 2.5~ms & 2.5--5~ms \\
놓이는 시간, 코일에 다이오드만 & 3--5~ms & 3--10~ms (추정) \\
놓이는 시간, 다이오드 + 제너 30~V & 1.5~ms & 1.5--4~ms \\
접점이 공중에 떠 있는 틈(먼저 끊고 나중에 잇기) & & 0.2--0.6~ms (추정) \\
\bottomrule
\end{tabular}
\end{table}

\General 신호등의 노란불은 ``최소 3초''로 정한다. 교차로가 넓으면 늘이고, 짧게 하는 것만 금지한다. 엘리베이터의
문이 ``적어도 몇 초는 열려 있어야 한다''는 것도 같다. 둘 다 사람이 얼마나 빨리 움직일지 정확히 모르므로
가장 느린 사람에 맞춘 하한을 둔다.

\Ours 기어식은 빠른 드럼 60~rpm에서 한 위치가 38.46~ms이고(자리올림 칸까지 넣은 평균은 위치당 약 44~ms,
표~\ref{tab:f-speed}), 릴레이식은 순서기가 한 위치를 63~ms로 낸다.
두 기계 모두 같은 테스트 레지스터(공통부)를 쓰므로, 레지스터가 요구하는 하한은 두 기계가 함께 지킨다.

\section{테스트 레지스터: 붙잡아 두고 나중에 판정한다}\label{sec:t-register}

테스트 레지스터는 봄베의 판단 장치다. \ref{sec:f-register}절에서 본 대로, 테스트 글자 케이블의 26선이
\emph{전부} 통전되면 그 위치는 버리고, 하나라도 통전되지 않으면 정지한다. 26선마다 \term{감지 릴레이}{sense
relay} $K$가 하나씩 있다.

\Why 처음 회로는 판정을 시험 전압이 걸려 있는 동안 끝내려 했다. 시험 전압이 걸리면 감지 릴레이가 붙고, 26개가
모두 붙으면 직렬로 이은 접점 사슬이 닫혀 \term{통과 릴레이}{pass relay} $P$가 붙고, 판정 순간에 $P$가 놓여
있으면 \term{정지 릴레이}{stop relay} $S$가 붙는다. 릴레이 세 단이 차례로 움직여야 한다. 최악의 시간으로 더하면
$5 + 5 + 5 = 15$~ms인데, 기어식에서 드럼 접점이 패드 위에 확실히 얹혀 있는 창은 6~ms다. 대표값으로도
8~ms라 들어가지 않는다. 기어를 늦추면 들어가지만, 최악값을 담으려면 속도가 3분의 1이 된다.

해법은 \term{표본 유지}{sample-and-hold}다. 시험 전압이 걸린 짧은 창 안에서는 감지 릴레이가 \emph{붙기만} 하고,
붙은 릴레이는 자기 유지(\ref{sec:relay-latch}절)로 그 결과를 붙든다. 판정은 창이 닫힌 \emph{뒤에} 느긋하게 한다.
창 안에서 끝나야 하는 것은 릴레이 하나의 동작(최악 5~ms)뿐이다. 판정 순간에 감지 릴레이들은 이미 멈춰 있는
값이므로 빠르기 경쟁도 없다.

\General 디지털카메라는 셔터가 열린 짧은 순간에 빛을 받아 두고, 사진 처리는 셔터가 닫힌 뒤에 한다. 아날로그-디지털
변환기 앞의 표본 유지 회로도 같다. 변하는 입력을 한 순간 커패시터에 붙잡아 두고, 변환은 그 붙잡힌 값으로
천천히 한다.

\Ours 그래서 뱅크(기어식 구동부나 릴레이식 순서기, D부와 E부)는 테스트 레지스터에 세 신호를 차례로 보낸다.

\begin{itemize}
\item \code{STROBE}: 시험 전압을 메뉴 배선망에 건다. 통전된 선의 $K$가 붙어 스스로를 유지한다.
\item \code{JUDGE}: 판정 경로를 연다. 그때 $P$가 놓여 있으면(26선이 전부는 아니면) $S$가 붙어 정지를 붙든다.
\item \code{CLEAR}: 유지 버스 $H$를 끊어 $K$와 $P$를 놓는다. 다음 위치를 위한 원상 복귀다.
\end{itemize}

세 신호는 ``언제''만 알린다. ``정지인가''는 레지스터 안의 릴레이 접점이 정한다. 그림~\ref{fig:t-cycle}은 한
위치 동안 일어나는 일이다.

\begin{figure}[ht]
\centering
\begin{tikztimingtable}[timing/slope=0.05,timing/coldist=4pt,xscale=0.36,yscale=1.15,semithick]
  \code{STROBE} (접점) & 6H 32.5L \\
  $K$ (통전된 선) & 2.5L 28.5H 7.5L \\
  $P$ (26선 전부일 때) & 7L 23H 8.5L \\
  \code{JUDGE} (접점) & 12L 7H 19.5L \\
  $S$ (정지일 때) & 16L 22.5H \\
  \code{CLEAR} (접점) & 21L 12H 5.5L \\
\extracode
  \begin{pgfonlayer}{background}
    \foreach \x in {0,6,12,19,21,33,38.5} \draw[cgray!50,dashed] (\x,1.2) -- (\x,-11.5);
  \end{pgfonlayer}
  \foreach \x in {0,6,12,19,21,33,38.5} \node[font=\tiny,text=cgray] at (\x,-12.0) {\x};
  \node[note,anchor=west] at (39.2,-12.0) {ms};
  \node[note] at (3,1.7) {감지};
  \node[note] at (9,1.7) {정착};
  \node[note] at (15.5,1.7) {판정};
  \node[note] at (27,1.7) {복귀};
\end{tikztimingtable}
\caption{테스트 레지스터의 한 위치(기어식 60~rpm, 38.5~ms, 접점이 실제로 닫힌 구간). $P$와 $S$는 둘 가운데
하나만 일어난다. 26선이 전부 통전되면 $P$가 붙어 판정 경로를 막고, 아니면 $P$가 놓인 채라 $S$가 붙는다.
$S$가 붙으면 \code{CLEAR}가 와도 $K$가 지워지지 않는다. 그림의 $K$는 정지하지 않은 경우다.}
\label{fig:t-cycle}
\end{figure}

그림~\ref{fig:t-sh}는 이 동작의 핵심만 추린 회로다. 세 가지를 눈여겨볼 만하다.

\begin{figure}[ht]
\centering
\begin{circuitikz}[x=1cm,y=1cm]
  \ctikzset{bipoles/length=9mm}
  \draw (0,3) node[left,font=\footnotesize] {$+24$ V} to[nos,l={\scriptsize \code{STROBE}}] (1.6,3) -- (2.0,3);
  \node[cloud,minimum width=26mm,minimum height=8mm,anchor=west] (net) at (2.0,3) {메뉴 배선망};
  \draw (net.east) -- (5.6,3) node[above,font=\footnotesize] {선 $L_A$} -- (6.4,3);
  \draw (6.4,3) to[D] (6.4,1.6) coordinate (n);
  \node[rcoil,anchor=north] (K) at (6.4,1.3) {$K_A$};
  \draw (n) -- (K.north); \draw (K.south) -- (6.4,0.0) node[below,font=\footnotesize] {0 V};
  \draw (0,1.6) node[left,font=\footnotesize] {$+24$ V} to[ncs,l={\scriptsize \code{CLEAR}의 $r$}] (2.2,1.6)
        -- (3.0,1.6) node[above,font=\footnotesize,text=cgreen] {$H$} -- (3.8,1.6) to[nos,l={\scriptsize $k_A$ (자기 유지)}] (5.6,1.6) -- (n);
  \node[circ] at (n) {};
  \node[note,anchor=west,align=left] at (7.3,2.4) {차단 다이오드: 유지 전압이\\배선망으로 거꾸로 가지 않는다};
  \node[note,anchor=west,align=left] at (7.3,0.6) {$K_A$가 붙으면 $k_A$가 닫혀\\$H$에서 스스로를 먹인다};
\end{circuitikz}
\caption{감지 릴레이 하나의 표본 유지(26개 가운데 하나). 실제 회로에는 이 밖에 접점 사슬, 통과 릴레이 $P$,
정지 릴레이 $S$, 해제 릴레이가 있다(그림~\ref{fig:t-register-full}).}
\label{fig:t-sh}
\end{figure}

\begin{itemize}
\item \textbf{유지 버스 $H$.} 26개 감지 릴레이의 자기 유지 접점이 모두 같은 버스 $H$에서 전원을 받는다. $H$의
공급은 \code{CLEAR} 릴레이의 NC 접점을 지난다. \code{CLEAR} 동안 그 접점이 열리면 26개가 한꺼번에 놓인다.
정지 중에는 $S$의 접점이 이 NC 접점을 우회해 $H$를 살려 두므로, 정지 순간의 26선 상태가 램프에 남는다.
\item \textbf{차단 다이오드.} 감지 릴레이가 자기 유지에 들면 코일 쪽 마디에 +24~V가 계속 걸린다. 이 마디가
선에 그대로 이어져 있으면 그 전압이 메뉴 배선망으로 거꾸로 들어가 다른 선을 통전시키고, \code{STROBE}가
끝난 뒤에도 드럼 접점에 전류가 흐른다. 선마다 다이오드 하나가 이 길을 막는다. \ref{sec:relay-logic}절의 샛길
문제를 다이오드로 푼 예다.
\item \textbf{세 신호의 역할 분담.} \code{STROBE}, \code{JUDGE}, \code{CLEAR}는 레지스터 안의 타이밍 릴레이 셋을
움직인다. 시험 전류를 잇고 끊는 것도, 판정 경로를 여는 것도, $H$를 끊는 것도 레지스터 안의 접점이다.
\end{itemize}

\begin{lesson}{계산으로 정한 회로를 시뮬레이션이 고쳤다}
결정 0001은 처음에 계산으로만 이 회로를 정했다. 보드를 회로도로 옮기고 넷리스트를 시간 있는 스위치 수준
시뮬레이션(\ref{ch:boardverify}장)으로 돌리자 틀린 곳 넷이 나왔다. 그 가운데 둘은 \emph{놓친 정지}(위험한 쪽)로
이어지는 것이었다. 하나는 접점 사슬을 +24~V에서 먹인 것이다. 26선이 전부 붙은 위치에서는 \code{CLEAR} 동안에도
$P$의 코일이 사슬로 전압을 받아, $P$가 첫 $K$가 놓인 뒤에야 놓이기 시작했다. 그래서 다음 위치까지 $P$가 남아
그 위치를 통과로 만들었다. 사슬을 $H$에서 먹이게 고치자 $K$와 $P$가 같은 순간에 놓이기 시작했다. 다른 하나는
해제 버튼이 $S$만 풀던 것이다. 정지 중에 그 위치의 \code{CLEAR}가 이미 지나갔으면, 재출발 뒤 첫 \code{STROBE}가
지워지지 않은 $K$ 위에 덧붙어 다음 위치의 정지를 놓친다. 해제 릴레이 하나를 더해 해제 때 $H$까지 끊게 했다.
\end{lesson}

보드의 실제 회로는 그림~\ref{fig:t-register-full}이다(보고서와 같은 그림). 시뮬레이션은 릴레이가 가장 느릴 때
레지스터가 견디는 최소 폭도 찾았다. 접점 면에서 \code{STROBE} 5~ms(감지 릴레이의 최악 동작), \code{JUDGE}
5~ms(정지 릴레이의 최악 동작, 모자라면 정지를 놓친다), \code{CLEAR} 10~ms(감지 릴레이의 최악 복귀, 모자라면
$H$가 살아날 때 $K$가 다시 붙어 지워지지 않는다)다.

\begin{figure}[p]
\centering
\resizebox{\linewidth}{!}{% 테스트 레지스터와 정지 릴레이의 회로도 (06-hardware.tex, 그림 fig:hw-register)
% 코일은 대문자(K_A, P, S), 그 릴레이의 접점은 소문자(k_A, p, s)로 적는다.
\begin{circuitikz}[x=1cm,y=1cm]
  \ctikzset{bipoles/length=9mm}
  % ---- (1) 시험 전류: +24 V, STROBE 접점, 메뉴 배선망, 감지 릴레이 코일 26개 ----
  \node[font=\footnotesize,anchor=east] (v1) at (0,6.6) {$+24$ V};
  \draw (v1.east) to[nos,l={\footnotesize \code{STROBE}}] (2.3,6.6) -- (2.8,6.6);
  \node[cloud,minimum width=68mm,minimum height=9mm,anchor=west] (net) at (2.8,6.6)
    {메뉴 배선망 (스크램블러, 대각선 보드)};
  \foreach \x/\l in {3.5/A,5.0/B,6.5/C,8.9/Z} {
    \node[rcoil] (K\l) at (\x,5.2) {$K_{\l}$};
    \draw (\x,6.15) -- (K\l.north);
    \draw (K\l.south) -- (\x,4.3);
  }
  \node at (7.7,5.2) {$\cdots$};
  \draw (2.9,4.3) -- (9.6,4.3) node[right,font=\footnotesize] {0 V};
  \node[note,anchor=west] at (10.3,6.0) {입력 케이블의 26선마다\\감지 릴레이 코일 하나};
  \node[font=\footnotesize\bfseries,anchor=west] at (-1.1,7.4) {(1) 시험 전류의 경로};
  % ---- (2) 접점 사슬과 통과 릴레이 ----
  \node[font=\footnotesize\bfseries,anchor=west] at (-1.1,3.5) {(2) 접점 사슬: 26선이 전부 통전되었는가};
  \node[font=\footnotesize,anchor=east] (v2) at (0,2.6) {$+24$ V};
  \draw (v2.east) -- (2.6,2.6) to[nos,l={$k_A$}] (4.1,2.6) to[nos,l={$k_B$}] (5.6,2.6)
        to[nos,l={$k_C$}] (7.1,2.6);
  \draw[dashed] (7.1,2.6) -- (8.1,2.6);
  \draw (8.1,2.6) to[nos,l={$k_Z$}] (9.6,2.6) -- (10.6,2.6);
  \node[rcoil,anchor=west] (P) at (10.6,2.6) {$P$};
  \draw (P.east) -- (12.3,2.6) node[right,font=\footnotesize] {0 V};
  \node[note,anchor=north] at (11.05,2.25) {통과 릴레이};
  % ---- (3) 정지 릴레이와 자기 유지 ----
  \node[font=\footnotesize\bfseries,anchor=west] at (-1.1,1.55) {(3) 정지 릴레이};
  \node[font=\footnotesize,anchor=east] (v3) at (0,0.7) {$+24$ V};
  \draw (v3.east) -- (0.8,0.7) coordinate (a) -- (2.6,0.7) to[nos,l={$j$}] (4.1,0.7)
        to[ncs,l={$p$}] (5.6,0.7) -- (8.6,0.7) coordinate (b) -- (10.6,0.7);
  \node[rcoil,anchor=west] (S) at (10.6,0.7) {$S$};
  \draw (S.east) -- (12.3,0.7) node[right,font=\footnotesize] {0 V};
  \node[note,anchor=north] at (11.05,0.35) {정지 릴레이};
  \draw (a) -- (0.8,-0.7) -- (2.6,-0.7) to[ncs,l_={원위치 버튼}] (4.1,-0.7) -- (5.9,-0.7)
        to[nos,l_={$s$ (자기 유지)}] (7.4,-0.7) -- (8.6,-0.7) -- (b);
  \node[circ] at (a) {};
  \node[circ] at (b) {};
  % ---- (4) RUN ----
  \node[font=\footnotesize\bfseries,anchor=west] at (-1.1,-1.75) {(4) 구동 허가};
  \node[font=\footnotesize,anchor=east] (v4) at (0,-2.5) {$+24$ V};
  \draw (v4.east) -- (2.6,-2.5) to[nos,l_={시작}] (4.1,-2.5) to[ncs,l_={$s$}] (5.6,-2.5) -- (6.6,-2.5)
        node[right,font=\footnotesize] {뱅크의 \code{RUN} (정지 릴레이가 동작하면 끊긴다)};
\end{circuitikz}
}
\caption{테스트 레지스터 보드의 회로(보고서의 그림과 같다). 코일은 대문자($K_A$, $P$, $S$), 그 릴레이의
접점은 소문자다. $t$, $j$, $r$은 \code{STROBE}, \code{JUDGE}, \code{CLEAR}가 움직이는 타이밍 릴레이의 접점,
$x$는 해제 릴레이의 접점이다.}
\label{fig:t-register-full}
\end{figure}

\section{두 면: 커넥터의 전압과 접점의 실제}\label{sec:t-faces}

\Why 레지스터가 요구하는 시간은 \emph{레지스터 안의 접점}이 닫혀 있는 구간에 대한 것이다. 그런데 뱅크가 낼 수
있는 것은 \emph{커넥터 핀의 전압}뿐이다. 둘 사이에 타이밍 릴레이가 하나 있고, 그 릴레이는 핀이 오른 뒤
2.5--4~ms에 붙고(켜짐 지연 $d_{\mathrm{on}}$), 핀이 내린 뒤 1.5--4~ms에 놓인다(꺼짐 지연 $d_{\mathrm{off}}$).
결정 0003은 이 둘을 \term{접점 면}{contact face}과 \term{커넥터 면}{connector face}이라는 이름으로 나눠, 회로의
조건은 접점 면에서 적고 뱅크가 지킬 규칙은 커넥터 면에서 적었다(그림~\ref{fig:t-faces}).

\begin{figure}[ht]
\centering
\begin{tikzpicture}[x=0.5cm,y=1cm]
  \draw[->] (-0.5,0) -- (22,0) node[right,font=\footnotesize] {시간};
  % connector face pulse
  \draw[line width=1pt,cblue] (0,1.6) -- (2,1.6) -- (2,2.4) -- (12,2.4) -- (12,1.6) -- (20,1.6);
  \node[font=\footnotesize,anchor=east,text=cblue] at (-0.6,2.0) {커넥터 면 (핀)};
  % contact face pulse (range)
  \fill[cbluebg] (6,0.4) rectangle (10,1.0);
  \fill[cyellowbg] (12,0.4) rectangle (16,1.0);
  \draw[line width=1pt,corange] (0,0.4) -- (6,0.4) -- (6,1.0) -- (14,1.0) -- (14,0.4) -- (20,0.4);
  \node[font=\footnotesize,anchor=east,text=corange] at (-0.6,0.7) {접점 면 (레지스터 안)};
  \draw[<->] (2,1.3) -- node[above,font=\scriptsize] {$d_{\mathrm{on}}$ 2.5--4} (6,1.3);
  \draw[<->] (12,1.3) -- node[above,font=\scriptsize] {$d_{\mathrm{off}}$ 1.5--4} (14,1.3);
  \node[note,anchor=north,align=center] at (8,0.3) {파란 띠: 접점이 붙는\\때가 흩어지는 범위};
  \node[note,anchor=north,align=center] at (14,0.3) {노란 띠: 놓이는\\때의 범위};
\end{tikzpicture}
\caption{두 면. 커넥터 핀의 펄스(위)는 타이밍 릴레이를 거쳐 접점의 펄스(아래)가 된다. 접점 면의 폭은 커넥터
면의 폭 $- d_{\mathrm{on}} + d_{\mathrm{off}}$이므로 최악의 경우($d_{\mathrm{on}}=4$, $d_{\mathrm{off}}=1.5$) 2.5~ms 줄어든다.}
\label{fig:t-faces}
\end{figure}

그러면 커넥터 면의 규칙은 접점 면의 조건에서 \emph{계산으로} 나온다. 예를 들어 보자.

\begin{itemize}
\item 접점 면에서 \code{STROBE}는 5~ms 이상 닫혀야 한다(C1). 폭은 최악 2.5~ms 줄므로, 핀에서는 $5 + 2.5 = 7.5$~ms
이상이어야 한다(L1). \code{JUDGE}(L3)도 같은 이유로 7.5, \code{CLEAR}(L5)는 $10 + 2.5 = 12.5$다.
\item 접점 면에서 \code{JUDGE}는 \code{STROBE}가 닫힌 뒤 10~ms(감지 5 + 통과 5) 이후에 닫혀야 한다(C2).
\code{STROBE}의 접점은 핀이 오른 뒤 \emph{늦으면} 4~ms에 닫히고, \code{JUDGE}의 접점은 핀이 오른 뒤
\emph{빠르면} 2.5~ms에 닫힌다. 가장 나쁜 짝을 맞추면 \code{JUDGE}의 핀은 \code{STROBE}의 핀보다 $4 + 10 - 2.5 =
11.5$~ms 뒤에 올라야 한다(L2).
\item \code{CLEAR}의 접점은 \code{JUDGE}의 접점이 열린 뒤에 열려야 한다(C4). \code{JUDGE}의 접점은 핀이 내린 뒤
늦으면 4~ms에 열리고, \code{CLEAR}의 접점은 핀이 오른 뒤 빠르면 2.5~ms에 열린다. 그래서 $4 - 2.5 = 1.5$~ms(L4).
\end{itemize}

이렇게 얻은 규칙이 표~\ref{tab:t-rules}의 L1--L9다. 모두 하한이거나 순서이고, 레지스터의 실제 지연을 몰라도
성립한다. 하한을 모두 더한 가장 짧은 주기는 $0 \to 11.5 \to 19 \to 20.5 \to 33 \to 34.5$, 곧 34.5~ms다.

\begin{table}[ht]
\centering
\caption{뱅크가 지킬 규칙(결정 0003, 커넥터 면, ms).}
\label{tab:t-rules}
\small
\begin{tabularx}{\linewidth}{>{\raggedright\arraybackslash}p{9mm}>{\raggedright\arraybackslash}X>{\raggedright\arraybackslash}p{38mm}}
\toprule
& 요구 & 어기면 \\
\midrule
L1 & \code{STROBE}의 폭 $\ge 7.5$ & 감지 릴레이가 못 붙음(거짓 정지) \\
L2 & \code{JUDGE}의 오름 $\ge$ \code{STROBE}의 오름 $+ 11.5$ & $P$가 붙기 전에 판정(거짓 정지) \\
L3 & \code{JUDGE}의 폭 $\ge 7.5$ & $S$가 못 붙음(\textbf{놓친 정지}) \\
L4 & \code{CLEAR}의 오름 $\ge$ \code{JUDGE}의 내림 $+ 1.5$ & 판정 중에 $P$가 놓임 \\
L5 & \code{CLEAR}의 폭 $\ge 12.5$ & $K$가 지워지지 않음(\textbf{놓친 정지}) \\
L6 & 다음 \code{STROBE}의 오름 $\ge$ \code{CLEAR}의 내림 $+ 1.5$ & 유지 버스가 없는 동안 붙은 $K$를 잃음 \\
L7 & \code{STROBE}를 낸 위치에는 \code{JUDGE}와 \code{CLEAR}를 반드시 낸다 & 판정 없이 지나간 위치 \\
L8 & 전원 투입, \code{HOME}, 해제 뒤에는 \code{STROBE} 전에 \code{JUDGE}를 내지 않는다 & 빈 레지스터를 판정 = 정지 \\
L9 & 시프터·드럼의 전환은 \code{CLEAR}의 오름부터 다음 \code{STROBE}의 오름 사이에만 & 전류가 흐르는 접점을 옮김 \\
\bottomrule
\end{tabularx}
\end{table}

L9는 \term{무전류 전환}{dry switching} 규칙이다. 전류가 흐르는 채로 접점을 떼면 불꽃이 튀어 접점이 상한다.
시험 전압이 없는 동안에만 드럼이나 시프터를 움직이라는 것이다.

\General 두 면의 구분은 디지털 회로의 \term{셋업 시간}{setup time}과 같은 생각이다. 칩의 데이터시트는 ``핀에서
클럭 에지 몇 ns 전에 데이터가 안정되어야 한다''고 적는다. 칩 안의 실제 플립플롭이 요구하는 조건을, 핀에서 그
플립플롭까지의 지연이 가장 나쁠 때를 넣어 핀의 조건으로 옮긴 것이다.

\Ours 결정 0003은 하나 더 열어 두었다. 물린 레지스터의 지연을 실제로 재서(\term{보정}{calibration}) 그만큼
앞당겨 내면 접점 면의 조건을 직접 지킬 수 있다. 기어식은 \code{STROBE}만 이렇게 보정하고(드럼 접점의 가운데에
접점 면의 창을 놓아야 하므로), 나머지는 보정 없이 넓게 낸다. 릴레이식 순서기는 셋 모두 보정 없이 지킨다.

\section{정리}

\begin{itemize}
\item 릴레이의 시간은 넓게 흩어지고 최소를 모른다. 그래서 시간 규칙은 모두 하한으로 적고, 기계가 느려지면 늘
더 잘 지켜지게 한다.
\item 테스트 레지스터는 표본 유지다. \code{STROBE} 동안 감지 릴레이가 붙어 스스로를 붙들고, \code{JUDGE}에서
판정하고, \code{CLEAR}에서 유지 버스를 끊어 지운다. 선마다 다이오드가 유지 전압의 역류를 막는다.
\item 회로의 조건은 접점 면에서, 뱅크의 규칙(L1--L9)은 커넥터 면에서 적는다. 둘 사이의 지연 범위를 가장 나쁘게
넣어 하나에서 다른 하나를 계산한다.
\item 이 규칙을 기어식은 슬롯 원판으로(\ref{sec:t-gear}절), 릴레이식은 순서기로(\ref{ch:relay-sequencer}장) 지킨다.
\end{itemize}
